\documentclass{article}
\usepackage[UTF8]{ctex}
\usepackage{amsmath,amssymb,amsfonts}
\usepackage{enumitem}
\usepackage{geometry}
\geometry{a4paper, margin=1in}
% === 错题本排版优化（不参与题目标识） ===
\setlength{\parindent}{0pt}
\setlength{\parskip}{0.35em plus 0.10em minus 0.08em}
\setlength{\emergencystretch}{3em}
\setlist{itemsep=0.25em, topsep=0.40em, parsep=0pt}
\allowdisplaybreaks[2]
% === 错题本排版优化结束 ===
\begin{document}
\section*{华南师范大学2024年数学分析真题}
\section*{159．华南师范大学2024年数学分析试卷}

\subsection{题目1}
1． \(\lim _{x \rightarrow 0} \frac{\sqrt{1+\tan x}-\sqrt{1-\tan x}}{\sin x}\) ．

\subsection{题目2}
2． \(\lim _{x \rightarrow 1}\left(\frac{k}{1-x^{k}}-\frac{l}{1-x^{l}}\right)\) ，其中 \(k, l\) 为正整数。

\subsection{题目3}
3． \(\lim _{n \rightarrow \infty}\left[\cos \pi \sqrt{n^{2}+1}\right]^{2}\) ．

\subsection{题目4}
4． \(\lim _{n \rightarrow \infty} \frac{n}{\sqrt[n]{n!}}\) ．

\subsection{题目5}
5． \(\lim _{\substack{x \rightarrow+\infty \\ y \rightarrow+\infty}}\left(\frac{x y}{x^{2}+y^{2}}\right)^{x^{2}}\) ．

\subsection{题目6}
1．已知 \(f(x)\) 在原点处存在二阶导数，且 \(\lim _{x \rightarrow 0}\left(\frac{\sin x}{x^{3}}+\frac{f(x)}{x^{2}}\right)=0\) ，求 \(f^{\prime}(0), f^{\prime \prime}(0)\) ．

\subsection{题目7}
2．计算 \(\int_{-2}^{-\sqrt{2}} \frac{1}{x \sqrt{x^{2}-1}} \mathrm{~d} x\) ．

\subsection{题目8}
3．计算 \(\iiint_{\Omega} x \mathrm{~d} x \mathrm{~d} y \mathrm{~d} z, \Omega\) 是由平面 \(x+2 y+z=1, x=0, y=0, z=0\) 围成的区域．

\subsection{题目9}
4．已知 \(u\) 是由 \(u=f(x, y, z, t), g(y, z, t)=0, h(z, t)=0\) 确定的关于 \(x, y\) 的函数，求 \(\frac{\partial u}{\partial x}, \frac{\partial u}{\partial y}\) ．

三．（10 分）已知 \(f(x)\) 是 \((0,+\infty)\) 上的单调递增函数，且存在正数列 \(\left\{a_{n}\right\}\) 满足

\[
\lim _{n \rightarrow \infty} a_{n}=+\infty, \lim _{n \rightarrow \infty} f\left(a_{n}\right)=A
\]

证明： \(\lim _{x \rightarrow+\infty} f(x)=A\) ．

四．（10 分）已知 \(f(x)\) 在 \([a, b]\) 上连续且单调递增，证明： \(\int_{a}^{b} x f(x) \mathrm{d} x \geq \frac{a+b}{2} \int_{a}^{b} f(x) \mathrm{d} x\) ．

五．（10 分）证明函数项级数 \(\sum_{n=1}^{\infty} \frac{n x}{1+n^{5} x^{2}}\) 在 \([0,+\infty)\) 上一致收敛。

六．（10 分）已知 \(f(x)\) 在 \([0,+\infty)\) 上连续，无穷积分 \(\int_{0}^{+\infty} f(x) \mathrm{d} x\) 收敛，证明：

\[
\lim _{u \rightarrow 0^{+}} \int_{0}^{+\infty} e^{-u x} f(x) \mathrm{d} x=\int_{0}^{+\infty} f(x) \mathrm{d} x
\]

七．（20 分）已知

\[
f(x, y)= \begin{cases}\left(x^{2}+y^{2}\right) \sin \frac{1}{x^{2}+y^{2}}, & (x, y) \neq(0,0) ; \\ 0, & (x, y)=(0,0) .\end{cases}
\]

证明：

\begin{itemize}
  \item[（1）] \(f(x, y)\) 在 \((0,0)\) 处连续．
  \item[（2）] \(f(x, y)\) 在 \((0,0)\) 处存在偏导数．
  \item[（3）] \(f(x, y)\) 在 \((0,0)\) 处偏导数不连续．
  \item[（4）] \(f(x, y)\) 在 \((0,0)\) 处可微．
\end{itemize}
\end{document}
